\section{Inciso a)}

\textit{Identifica el patrón de diseño más adecuado, justifica tu selección basándote en las
características del problema y explica la propuesta de solución considerando las clases que
intervienen y las relaciones entre ellas. ¿Por qué no usar Strategy o polimorfismo simple?}

\vspace{0.4cm}

\subsection*{Patrón seleccionado}

El patrón adecuado es \textbf{Visitor} (Visitante), patrón de comportamiento del catálogo
\textit{Gang of Four}. Su intención declarada es: \textit{representar una operación que se
ejecuta sobre los elementos de una estructura de objetos, permitiendo definir nuevas operaciones
sin modificar las clases de los elementos sobre los que opera}. Esa intención es, palabra por
palabra, el requisito central del enunciado.

\subsection*{Justificación a partir de las características del problema}

El problema presenta dos ejes de variación independientes, y la decisión de diseño depende por
completo de la velocidad relativa con que crece cada uno:

\begin{table}[h!]
    \centering
    \begin{tabular}{lp{6.6cm}l}
        \toprule
        \textbf{Eje} & \textbf{Elementos} & \textbf{Evolución} \\
        \midrule
        Tipos de elemento & CuentaAhorro, CuentaCorriente, CuentaInversión, CuentaEmpresarial
                          & Estable \\
        Operaciones       & Ejecutar, Revertir, Auditar, congelar cuenta, calcular intereses, \ldots
                          & Crece \\
        \bottomrule
    \end{tabular}
\end{table}

El enunciado fija el eje de los tipos de cuenta (son cuatro y no se anuncian más) y declara
explícitamente abierto el eje de las operaciones (``debe permitir añadir nuevas operaciones
sin modificar las clases de las cuentas ni de las transacciones existentes''). Esto es el
\textbf{principio Open/Closed} enunciado sobre el eje de las operaciones: el sistema debe estar
abierto a la extensión por operaciones y cerrado a la modificación de los elementos.

Se trata de una instancia del \textit{expression problem}: con herencia convencional resulta
barato añadir \emph{tipos} nuevos y caro añadir \emph{operaciones} nuevas; Visitor invierte
exactamente ese balance. Expresado con la métrica de ``número de cambios'' que el problema
propone, y siendo $T$ el número de tipos de cuenta y $V$ el número de operaciones:

\begin{itemize}[leftmargin=1.2cm]
    \item \textbf{Operación nueva con Visitor:} 1 clase nueva, 0 modificaciones.
          Con métodos en cada elemento: $T$ clases modificadas más la interfaz, es decir
          $T+1$ modificaciones.
    \item \textbf{Tipo de cuenta nuevo con Visitor:} 1 clase nueva más $V+1$ modificaciones
          (todos los visitantes y su interfaz). Con métodos en cada elemento: 1 clase nueva.
\end{itemize}

Visitor es por tanto la elección correcta \emph{porque} en este dominio se cumple $V \gg T$ y
$V$ crece mientras $T$ permanece fijo. Si la relación fuera la inversa, el patrón sería una
mala decisión: la justificación no es el patrón en abstracto, sino el perfil de evolución del
sistema.

A esto se añade el \textbf{principio de responsabilidad única}: si las operaciones vivieran
dentro de las cuentas, la clase CuentaAhorro acabaría conociendo las reglas de auditoría, el
cálculo de intereses, las políticas de congelación y la emisión de reportes. Con Visitor, cada
operación queda en una sola clase cohesiva, legible, testeable y versionable como unidad.

\subsection*{Mecanismo: el doble despacho}

El punto técnico que justifica el patrón es el siguiente: el comportamiento requerido no
depende de un tipo, sino de un \textbf{par} $(\textit{tipo de cuenta},\ \textit{tipo de
operación})$. Ejecutar sobre una cuenta de inversión no es lo mismo que ejecutar sobre una de
ahorro, ni que auditar sobre una de inversión. Se necesita entonces resolver en tiempo de
ejecución los tipos dinámicos de \emph{dos} participantes: \textbf{despacho múltiple}.

Los lenguajes orientados a objetos de uso corriente (Java, C\#, C++, Python, TypeScript) solo
ofrecen \textbf{despacho simple}: el método ejecutado depende únicamente del tipo dinámico del
receptor, mientras los parámetros se resuelven estáticamente. Visitor simula el despacho doble
encadenando dos despachos simples:

\begin{verbatim}
interface ElementoContable {
    <R> R accept(VisitanteOperacion<R> v);
}

final class CuentaAhorro implements ElementoContable {
    public <R> R accept(VisitanteOperacion<R> v) {
        return v.visitAhorro(this);   // 2do despacho
    }
}

interface VisitanteOperacion<R> {
    R visitAhorro(CuentaAhorro c);
    R visitCorriente(CuentaCorriente c);
    R visitInversion(CuentaInversion c);
    R visitEmpresarial(CuentaEmpresarial c);
}
\end{verbatim}

La llamada \texttt{cuenta.accept(visitante)} resuelve en el primer despacho el tipo concreto de
la cuenta; ya dentro de \texttt{CuentaAhorro.accept}, el receptor \texttt{this} tiene tipo
estático CuentaAhorro, de modo que la selección de \texttt{visitAhorro} es inmediata y el
segundo despacho resuelve el tipo concreto de la operación. El método \texttt{accept} no
implementa lógica de negocio: su única función es revelar el tipo del elemento.

De aquí se sigue un criterio de verificación del diseño: \textbf{si en la implementación
aparece un \texttt{switch} o un \texttt{instanceof} sobre tipos de cuenta, el patrón no está
aplicado}. Eliminar ese \textit{type switch} es precisamente lo que aporta Visitor.

\subsection*{Por qué no polimorfismo simple}

El polimorfismo simple (métodos \texttt{ejecutar()}, \texttt{revertir()}, \texttt{auditar()}
declarados en la jerarquía de cuentas) resuelve un solo eje, y es el eje equivocado:

\begin{itemize}[leftmargin=1.2cm]
    \item Cada operación nueva obliga a modificar las cuatro clases de cuenta y su interfaz
          común, violando Open/Closed justo donde el enunciado exige extensibilidad.
    \item Concentra en la cuenta responsabilidades heterogéneas (auditoría, intereses,
          congelación), violando la responsabilidad única.
    \item Dispersa cada operación entre $T$ archivos: la lógica de auditoría deja de existir
          como unidad y no puede leerse, probarse ni modificarse de forma aislada.
\end{itemize}

\subsection*{Por qué no Strategy}

Strategy también encapsula un algoritmo en un objeto, de ahí la confusión; pero su semántica es
distinta en cuatro aspectos decisivos:

\begin{enumerate}[leftmargin=1.2cm]
    \item \textbf{Strategy varía \emph{cómo} se hace algo; Visitor varía \emph{sobre qué tipo}
          se hace.} Strategy parametriza un contexto con algoritmos intercambiables de la misma
          firma, y cada algoritmo trata su entrada como un único tipo uniforme.
    \item \textbf{Strategy no ofrece doble despacho.} Al recibir una Cuenta genérica, cada
          estrategia tendría que preguntar por su clase concreta, reintroduciendo el
          \textit{type switch} que el requisito de extensibilidad obliga a eliminar.
    \item \textbf{En Strategy el cliente elige el algoritmo} como grado de libertad en tiempo
          de ejecución. Aquí no existe tal elección: el comportamiento de ``ejecutar'' sobre
          una cuenta de ahorro está determinado por el tipo de cuenta, no por una preferencia
          del cliente.
    \item \textbf{Strategy no recorre estructuras de objetos.} El patrón elegido debe
          integrarse con Composite e Iterator (incisos siguientes) para operar sobre grupos de
          cuentas y acumular resultados a lo largo de una jerarquía; Visitor lo hace de forma
          natural mediante la propagación del \texttt{accept}, mientras Strategy no contempla
          recorridos.
\end{enumerate}

\subsection*{Propuesta de solución: clases y relaciones}

\textbf{Jerarquía de elementos} (cerrada y estable):

\begin{itemize}[leftmargin=1.2cm]
    \item \texttt{ElementoContable}: interfaz con la única operación
          \texttt{accept(VisitanteOperacion)}.
    \item \texttt{CuentaAhorro}, \texttt{CuentaCorriente}, \texttt{CuentaInversion},
          \texttt{CuentaEmpresarial}: implementan \texttt{accept} con una sola línea que invoca
          el método de visita correspondiente.
    \item \texttt{GrupoCuentas}: nodo compuesto que implementa la misma interfaz y propaga el
          \texttt{accept} a sus hijos (desarrollado en el inciso c).
\end{itemize}

\textbf{Jerarquía de visitantes} (abierta, es la que crece):

\sloppy
\begin{itemize}[leftmargin=1.2cm]
    \item \texttt{VisitanteOperacion<R>}: interfaz con un método de visita por tipo concreto de
          cuenta y un tipo de retorno genérico \texttt{R}.
    \item \texttt{EjecutarTransaccionVisitor}, \texttt{RevertirTransaccionVisitor},
          \texttt{AuditarTransaccionVisitor}: las tres operaciones exigidas, cada una con
          comportamiento propio por tipo de cuenta.
    \item \texttt{CongelarCuentaVisitor}, \texttt{CalcularInteresesVisitor}: operaciones futuras
          que se incorporan como clases nuevas, sin modificar ninguna cuenta.
\end{itemize}
\fussy

\textbf{Ubicación del parámetro de la operación.} La \texttt{Transaccion} concreta (depósito,
retiro, transferencia, pago de servicios) no viaja en la firma de \texttt{accept}, sino en el
estado del visitante concreto, construido como
\texttt{new EjecutarTransaccionVisitor(transaccion)}. Así la interfaz de los elementos
permanece invariante ante cualquier operación nueva, que es la propiedad que se quiere
preservar. Como contrapartida, el visitante tiene estado mutable, y de ahí surge la restricción
de concurrencia tratada a continuación.

\textbf{Relaciones.} El patrón introduce una dependencia cíclica a nivel de tipos que le es
intrínseca: los elementos dependen de la interfaz \texttt{VisitanteOperacion} y los visitantes
dependen de los tipos concretos de elemento. El cliente conoce la estructura de cuentas y el
visitante concreto, e inicia el proceso con \texttt{elemento.accept(visitante)}; ninguna otra
clase conoce el emparejamiento entre tipo de cuenta y operación.

\subsection*{Restricción: concurrencia y estado mutable del visitante}

Conviene separar dos problemas que suelen confundirse.

\textbf{(1) Estado mutable del visitante.} Como el visitante transporta datos (la transacción,
acumuladores de auditoría), compartir una única instancia entre hilos produce una condición de
carrera. Soluciones, de preferible a tolerable:

\begin{itemize}[leftmargin=1.2cm]
    \item \textbf{Visitante sin estado que devuelve resultado}: los métodos de visita retornan
          un valor inmutable de tipo \texttt{R} en lugar de acumular en un campo. Un visitante
          puro es paralelizable sin sincronización y sus resultados parciales se combinan por
          reducción (esquema \textit{fork-join} o \textit{map-reduce}). Es la opción correcta.
    \item \textbf{Confinamiento por hilo}: una instancia de visitante por recorrido o por hilo,
          combinando los resultados al final.
    \item Si la acumulación compartida es inevitable, usar acumuladores concurrentes
          (\texttt{LongAdder}, colecciones concurrentes). Sincronizar el visitante completo
          serializaría los recorridos y anularía el paralelismo exigido.
\end{itemize}

\textbf{(2) Atomicidad y consistencia por cuenta.} La exclusión mutua no pertenece al
visitante, sino al elemento: la cuenta es la unidad de consistencia del dominio y es ella quien
protege sus invariantes. En consecuencia:

\begin{itemize}[leftmargin=1.2cm]
    \item \textbf{Bloqueo por cuenta} (\textit{lock striping}: un cerrojo por identificador de
          cuenta en un mapa concurrente) o \textbf{control optimista} con número de versión y
          reintento mediante comparación e intercambio. Con cualquiera de los dos, las
          operaciones sobre cuentas distintas progresan en paralelo y las que afectan a la
          misma cuenta se serializan, que es exactamente lo que pide el enunciado.
    \item \textbf{Transferencia entre dos cuentas:} al requerir dos cerrojos aparece riesgo de
          \textit{deadlock}. Se evita imponiendo un orden total de adquisición (por
          identificador ascendente) o un bloqueo en dos fases con tiempo límite.
    \item \textbf{Separación de responsabilidades:} para que el visitante no mezcle lógica de
          negocio con lógica transaccional, el control de la transacción puede introducirse como
          un \textbf{decorador} del visitante (\texttt{VisitanteTransaccional}), que abre la
          transacción, delega en el visitante envuelto y confirma o revierte según el resultado.
    \item \textbf{Orden FIFO por cuenta:} si además se exige orden de llegada, una cola por
          cuenta (modelo de actores o \textit{single writer} por partición) convierte la
          serialización en una propiedad estructural del diseño en lugar de depender de
          cerrojos.
\end{itemize}

\subsection*{Observación sobre la combinación con otros patrones}

El enunciado pide ejecutar, \emph{revertir} y auditar. El \textit{rollback} y la traza de
auditoría corresponden conceptualmente a \textbf{Command} (encapsular la petición con
\texttt{ejecutar}/\texttt{deshacer} y registrarla en un \textit{log}) y a \textbf{Memento}
(conservar el estado previo de la cuenta). Visitor es la respuesta al eje ``añadir operaciones
sin modificar los elementos''; en un diseño completo ambos coexisten: el visitante aporta el
comportamiento polimórfico según el tipo de cuenta, y el \textit{command} define la unidad
transaccional reversible y auditable.

\subsection*{Limitaciones del patrón elegido}

\begin{itemize}[leftmargin=1.2cm]
    \item Un tipo de cuenta nuevo obliga a modificar la interfaz de visitante y todas sus
          implementaciones. El patrón apuesta a que la jerarquía de elementos es cerrada.
    \item \texttt{accept} debilita el encapsulamiento: la cuenta debe exponer al visitante los
          datos que este necesita (saldo, movimientos), lo que puede ensanchar su interfaz
          pública.
    \item Añade indirección. Para un sistema con dos operaciones estables, Visitor sería
          sobreingeniería; aquí se justifica únicamente por el requisito explícito de
          extensibilidad en el eje de las operaciones.
\end{itemize}
